\chapter*{Résumés des sessions}
\addcontentsline{toc}{chapter}{Résumés des sessions}

Cette section rassemble un résumé rapide de chaque session. Elle permet de retrouver les événements essentiels de la campagne avant de reprendre l'aventure.

\section*{Résumé de la session d'initiation --- La Faille de Waterdeep}

Kael, Morrÿs et Obélik se réveillent dans une cellule de la prison de Waterdeep, sans savoir comment ils y sont arrivés. Profitant du chaos provoqué par une invasion de kobolds, ils s'évadent et rencontrent Alan, un tavernier lui aussi emprisonné.

En explorant la prison, les quatre fugitifs découvrent une immense faille dont le fond reste invisible, puis une fosse commune remplie de cadavres. Leur progression les conduit ensuite jusqu'à un escalier magique qui se transforme en toboggan et déclenche une alarme. Malgré le code trouvé sur un garde, ils choisissent de le gravir à l'aide d'une corde et de pitons.

Dans le hall, ils rejoignent le capitaine de la prison et l'aident à combattre une horde de kobolds. Une fois la bataille terminée, le capitaine comprend qu'ils sont des prisonniers évadés, mais décide de les laisser partir car la prison n'est plus sûre.

À l'extérieur, une créature invisible transperce le capitaine à plusieurs reprises. Un mystérieux halfelin abat finalement la créature d'une seule flèche, soigne le blessé et conseille au groupe de rejoindre Phandalin. Il repart ensuite vers la prison en marchant sur les murs puis au plafond, laissant les aventuriers prendre la route avec beaucoup plus de questions que de réponses.

\section*{Résumé de la session 2 --- La route de Phandalin}

Kael, Morrÿs, Obélik et Alan quittent Waterdeep en compagnie du capitaine blessé, qui leur révèle se nommer Virgile. En chemin, ils reparlent de la faille de la prison et apprennent qu'il avait tenté d'en mesurer la profondeur en y jetant une chope préalablement vidée.

Le groupe fait halte dans une taverne, malgré les guenilles de prisonniers que portent encore les fugitifs. Virgile leur demande de rester discrets pendant qu'il commande pour eux. Alan explique alors avoir été emprisonné pour avoir placardé des affiches publicitaires pour la taverne qu'il gère et espère un jour pouvoir racheter.

Virgile leur parle des Faiseurs de Veuves, une guilde installée à Phandalin qui pourrait soigner ses blessures. Malgré son nom inquiétant, l'organisation œuvre pour le bien du royaume ; son surnom lui vient de l'échec tragique de sa toute première mission.

À l'approche de Phandalin, le groupe aperçoit une colonne de fumée qu'Obélik attribue immédiatement à un barbecue. Devant une vaste demeure située à l'extérieur de la ville, une elfe membre des Faiseurs de Veuves les attend. Elle affirme être là depuis cinq jours alors que personne, dans le groupe, ne se souvient d'avoir annoncé leur arrivée.

\section*{Résumé de la session 3 --- Les Faiseurs de Veuves}

L'elfe conduit le groupe dans le quartier général des Faiseurs de Veuves, où ils rencontrent Yarno El Raïb et Destos. Yarno leur apprend que Corrin Thinspike, le mystérieux halfelin de la prison, les a recommandés. Il leur explique également que la guilde recherche et neutralise les objets magiques dangereux.

Comme les fugitifs ne peuvent ni retourner à Waterdeep ni rester longtemps à Phandalin, Yarno leur propose une mission à Chult. Le prince marchand Wakanga O'Tamu recherche un Gardien Bouclier et son amulette de contrôle. À Port Nyanzaru, Iolka et Nera Vesper, surnommées « la Paire », doivent leur servir de contacts.

Avant de les envoyer, Yarno leur impose une épreuve : toucher trois fois Lyra, l'elfe qui les a accueillis. Obélik, Destos et les autres parviennent à remplir l'objectif malgré plusieurs passages au sol. Morrÿs voit surtout son poignard dévié puis planté dans le plafond, où Obélik doit aller le récupérer.

Le groupe reçoit de l'or, des potions, une carte et une photographie de la Paire. Destos rejoint officiellement l'expédition, Alan trouve une place dans les cuisines et Virgile décide de retourner à Waterdeep pour ne pas être considéré comme un déserteur. Après un dernier repas préparé par Alan, le départ est fixé à l'aube.

\section*{Résumé de la session 4 --- L'Effleureuse}

Lyra accompagne Kael, Morrÿs, Obélik et Destos jusqu'à Leilon, où Corvis Quillfore les attend à bord de l'\textit{Effleureuse}. Faute d'équipage, les quatre aventuriers doivent participer aux manœuvres du trois-mâts magique pendant la traversée vers Chult.

Peu après leur départ, un navire de la garde de Waterdeep les intercepte. Les fugitifs se déguisent et tentent de tromper la Brigade d'Intervention Technique et Expéditive, mais leurs mensonges deviennent rapidement suspects. Morrÿs est reconnu et ligoté tandis que Destos détourne l'attention des gardes avec sa magie.

Obélik tranche la corde reliant les deux navires et le groupe prend la fuite. Kael projette le capitaine adverse par-dessus bord, Obélik expédie un soldat vers son propre bâtiment et Morrÿs se blesse avec son poignard en essayant de couper ses liens. La garde se lance alors à leur poursuite sous les tirs d'arbalète.

Corvis active finalement le pouvoir de l'\textit{Effleureuse}, qui bondit au-dessus de hauts-fonds tandis que le navire ennemi s'échoue. Le capitaine tombe à l'eau mais sa ligne de vie permet à Obélik de le ramener. Après quelques gifles médicinales et une recherche infructueuse de traqueur magique, le groupe reprend sa route vers Chult.

\section*{Résumé de la session 5 --- Port Nyanzaru}

Pendant la traversée, les aventuriers découvrent davantage l'\textit{Effleureuse} et organisent des quarts de nuit. Kael, Morrÿs, Obélik puis Destos font d'étranges cauchemars liés à leur passé, tous marqués par une même obscurité. Le voyage se poursuit néanmoins sans autre incident majeur, hormis le mal de mer d'Obélik et quelques tortues prises un instant pour un dragon des mers.

Après une vingtaine de jours, le groupe atteint Port Nyanzaru, unique cité de Chult. L'arrivée est marquée par la chaleur tropicale, les odeurs d'épices et une procession funèbre silencieuse. Corvis repart chercher un nouvel équipage tandis qu'un domestique conduit les aventuriers jusqu'à la villa de Wakanga O'Tamu.

Wakanga leur présente Syndra, une vieille elfe autrefois ressuscitée dont le corps dépérit désormais. Toutes les personnes revenues à la vie subissent le même mal et plus aucun prêtre ne parvient à ressusciter les morts. La source se trouverait dans la jungle et pourrait être liée au \textit{Soulmonger}, un artefact qui capturerait les âmes des défunts.

Le groupe accepte d'aider Syndra sans abandonner la recherche du Gardien Bouclier. Wakanga lui remet une carte, recommande de trouver un guide et révèle qu'Iolka et Nera n'ont plus donné de nouvelles depuis une vingtaine de jours. En quittant la villa pour rejoindre le quartier du marché, Kael ne remarque rien de précis, mais acquiert la certitude que quelqu'un les observe.
